% 08_features_demo.tex
\section{Features Demo}


\subsection{System Overview}

The system is designed as an interactive visualization tool that integrates user interface controls, algorithm execution, and real-time graphical rendering. A demonstration video of the system is available at:

\textbf{Demo Video:} link\_video

The video illustrates the core functionalities of the application, including user interaction, animation control, and visualization of multiple data structures.

\subsection{Core Features}

The system is organized into four main functional components:

\textbf{Navigation and Theme Interface.}
The application provides a main menu for selecting different data structures. It also supports theme switching between light and dark modes, along with the ability to show or hide the pseudocode panel.

\textbf{Control Panel.}
Users can input data through multiple methods, including random generation, console input, and file loading. The animation system includes controls such as play, pause, step forward, and step backward, as well as a speed adjustment slider ranging from 0.5x to 4.0x.

\textbf{Pseudocode Tracking.}
The system highlights corresponding lines of pseudocode during execution, enabling users to directly map algorithm steps to visual changes on the screen.

\textbf{Data Structure Operations Visualization.}
The application supports visualization for several data structures:
\begin{itemize}
    \item Singly Linked List: insertion, deletion, and node updates.
    \item Heap: heapify operations during insertion and removal.
    \item AVL Tree: self-balancing through rotation operations.
    \item Graph: shortest path algorithms such as Dijkstra and Bellman-Ford.
\end{itemize}

\subsection{Demonstration Scenarios}

The demonstration video presents several typical use cases:

\textbf{AVL Tree Balancing.}
When inserting nodes, the system automatically performs rotations to maintain balance. The animation highlights imbalance detection and the corresponding rotation steps.

\textbf{Heap Operations.}
The heap structure is dynamically adjusted through heapify processes, allowing users to observe how elements are rearranged to maintain heap properties.

\textbf{Graph Shortest Path Algorithms.}
Shortest path algorithms are visualized as propagation processes across the graph, where nodes and edges are progressively updated to reflect distance calculations.

\subsection{Discussion}

The integration of interactive controls, synchronized pseudocode, and animated visualization provides a comprehensive learning environment. The demonstration confirms that the system effectively bridges the gap between abstract algorithmic concepts and their practical execution.